% ========================================================
% BÀI TẬP TOÁN 5 - KẾT NỐI TRI THỨC VỚI CUỘC SỐNG
% Bài 58: Nhân, chia số đo thời gian với một số
% Định dạng: BTTeX Part III - Trắc nghiệm điền số ngắn
% ========================================================

% Câu 1
\begin{ex}
  Tính: $1\text{ giờ } 15\text{ phút} \times 3$. Kết quả có số phút là bao nhiêu?
  \shortans{45}
  \loigiai{
    $1\text{h} \times 3 = 3\text{h}$; $15 \times 3 = 45$ phút.
  }
\end{ex}

% Câu 2
\begin{ex}
  Tính: $2\text{ giờ } 25\text{ phút} \times 3$. Đổi kết quả thành mấy giờ $15$ phút?
  \shortans{7}
  \loigiai{
    $6\text{ giờ } 75\text{ phút} = 7\text{ giờ } 15\text{ phút}$.
  }
\end{ex}

% Câu 3
\begin{ex}
  Tính: $3\text{ phút } 12\text{ giây} \times 4$. Kết quả có số giây là bao nhiêu?
  \shortans{48}
  \loigiai{
    $12 \times 4 = 48$ giây.
  }
\end{ex}

% Câu 4
\begin{ex}
  Tính: $1\text{ ngày } 8\text{ giờ} \times 3$. Kết quả đúng bằng bao nhiêu ngày?
  \shortans{4}
  \loigiai{
    $3\text{ ngày } 24\text{ giờ} = 4\text{ ngày}$.
  }
\end{ex}

% Câu 5
\begin{ex}
  Một người thợ làm xong một chiếc ghế hết $1\text{ giờ } 45\text{ phút}$. Người đó làm $4$ chiếc ghế như thế hết bao nhiêu giờ?
  \shortans{7}
  \loigiai{
    $1\text{h}45 \times 4 = 4\text{h}180\text{ phút} = 7\text{ giờ}$.
  }
\end{ex}

% Câu 6
\begin{ex}
  Tính: $6\text{ giờ } 40\text{ phút} : 2$. Kết quả có số phút là bao nhiêu?
  \shortans{20}
  \loigiai{
    $40 : 2 = 20$ phút.
  }
\end{ex}

% Câu 7
\begin{ex}
  Tính: $7\text{ giờ } 30\text{ phút} : 3$. Kết quả là mấy giờ $30$ phút?
  \shortans{2}
  \loigiai{
    $7\text{h} : 3 = 2\text{h}$ dư 1h (60 phút); $90 : 3 = 30$ phút. Kết quả $2\text{h}30$.
  }
\end{ex}

% Câu 8
\begin{ex}
  Tính: $9\text{ phút } 36\text{ giây} : 4$. Kết quả có số giây là bao nhiêu?
  \shortans{24}
  \loigiai{
    $9\text{ phút} : 4 = 2\text{ phút}$ dư 1 phút (60 giây); $96 : 4 = 24$ giây.
  }
\end{ex}

% Câu 9
\begin{ex}
  Một tổ may $5$ chiếc áo hết $7\text{ giờ } 30\text{ phút}$. Trung bình may một chiếc áo hết bao nhiêu phút?
  \shortans{90}
  \loigiai{
    $7\text{ giờ } 30\text{ phút} = 450$ phút. $450 : 5 = 90$ phút.
  }
\end{ex}

% Câu 10
\begin{ex}
  Người ta đóng gói $6$ gói quà hết $1\text{ giờ } 12\text{ phút}$. Trung bình đóng mỗi gói hết bao nhiêu phút?
  \shortans{12}
  \loigiai{
    $72\text{ phút} : 6 = 12$ phút.
  }
\end{ex}

% Câu 11
\begin{ex}
  Một vận động viên hoàn thành $4$ vòng sân hết $8\text{ phút } 40\text{ giây}$. Trung bình mỗi vòng chạy hết bao nhiêu giây?
  \shortans{130}
  \loigiai{
    $8\text{ phút } 40\text{ giây} = 520$ giây. $520 : 4 = 130$ giây.
  }
\end{ex}

% Câu 12
\begin{ex}
  Tính: $2\text{ giờ } 15\text{ phút} \times 4$. Kết quả đúng bằng bao nhiêu giờ?
  \shortans{9}
  \loigiai{
    $8\text{ giờ } 60\text{ phút} = 9\text{ giờ}$.
  }
\end{ex}

% Câu 13
\begin{ex}
  Tính: $15\text{ giờ } 45\text{ phút} : 5$. Kết quả có số phút là bao nhiêu?
  \shortans{9}
  \loigiai{
    $45 : 5 = 9$ phút.
  }
\end{ex}

% Câu 14
\begin{ex}
  Bác An đóng $8$ cái bàn hết $20$ giờ. Hỏi trung bình đóng một cái bàn hết bao nhiêu phút?
  \shortans{150}
  \loigiai{
    $20\text{ giờ} = 1\,200$ phút; $1\,200 : 8 = 150$ phút.
  }
\end{ex}

% Câu 15
\begin{ex}
  Tính: $4\text{ ngày } 12\text{ giờ} \times 2$. Kết quả bằng bao nhiêu ngày?
  \shortans{9}
  \loigiai{
    $8\text{ ngày } 24\text{ giờ} = 9\text{ ngày}$.
  }
\end{ex}

% Câu 16
\begin{ex}
  Tính: $13\text{ giờ } 20\text{ phút} : 4$. Kết quả là mấy giờ $20$ phút?
  \shortans{3}
  \loigiai{
    $13\text{h} : 4 = 3\text{h}$ dư 1h (60 phút); $80 : 4 = 20$ phút. Kết quả là $3\text{h}20$.
  }
\end{ex}

% Câu 17
\begin{ex}
  Một máy in in $50$ trang sách hết $2\text{ phút } 30\text{ giây}$. Để in $100$ trang sách như thế hết bao nhiêu giây?
  \shortans{300}
  \loigiai{
    $150 \times 2 = 300$ giây.
  }
\end{ex}

% Câu 18
\begin{ex}
  Một máy bơm bơm $3$ bể nước như nhau hết $4\text{ giờ } 30\text{ phút}$. Bơm mỗi bể hết bao nhiêu phút?
  \shortans{90}
  \loigiai{
    $270 : 3 = 90$ phút.
  }
\end{ex}

% Câu 19
\begin{ex}
  Tính: $2\text{ phút } 45\text{ giây} \times 4$. Kết quả bằng bao nhiêu phút?
  \shortans{11}
  \loigiai{
    $8\text{ phút } 180\text{ giây} = 11\text{ phút}$.
  }
\end{ex}

% Câu 20
\begin{ex}
  Tính: $21\text{ phút } 15\text{ giây} : 3$. Kết quả có số giây là bao nhiêu?
  \shortans{5}
  \loigiai{
    $15 : 3 = 5$ giây.
  }
\end{ex}

% Câu 21
\begin{ex}
  Một ca nô đi $3$ chuyến hết $4\text{ giờ } 15\text{ phút}$. Hỏi ca nô đi $6$ chuyến như thế hết bao nhiêu phút?
  \shortans{510}
  \loigiai{
    $4\text{h}15 = 255$ phút. $255 \times 2 = 510$ phút.
  }
\end{ex}

% Câu 22
\begin{ex}
  Một người thợ làm việc $8$ giờ mỗi ngày, làm trong $5$ ngày được trả $2$ triệu đồng. Mỗi giờ làm việc được bao nhiêu nghìn đồng?
  \shortans{50}
  \loigiai{
    Tổng số giờ là $8 \times 5 = 40$ giờ. $2\,000 : 40 = 50$ nghìn đồng.
  }
\end{ex}

% Câu 23
\begin{ex}
  Tính: $(1\text{ giờ } 20\text{ phút} + 2\text{ giờ } 10\text{ phút}) \times 2$. Kết quả có số phút là bao nhiêu?
  \shortans{0}
  \loigiai{
    $3\text{h}30 \times 2 = 7\text{ giờ } 0\text{ phút}$. Số phút là 0.
  }
\end{ex}

% Câu 24
\begin{ex}
  Tính: $(5\text{ giờ } 30\text{ phút} - 1\text{ giờ } 50\text{ phút}) : 2$. Kết quả có số phút là bao nhiêu?
  \shortans{50}
  \loigiai{
    $3\text{h}40 : 2 = 1\text{ giờ } 50\text{ phút}$.
  }
\end{ex}

% Câu 25
\begin{ex}
  Bác Bình đi bộ $3$ vòng công viên hết $45$ phút. Nếu đi $5$ vòng như thế hết bao nhiêu phút?
  \shortans{75}
  \loigiai{
    Mỗi vòng hết $45 : 3 = 15$ phút. $5$ vòng hết $15 \times 5 = 75$ phút.
  }
\end{ex}
